\documentclass[11pt]{article}
\usepackage[margin=1in]{geometry}
\usepackage{amsmath,amssymb}
\usepackage{enumitem}
\usepackage[colorlinks=true,linkcolor=blue]{hyperref}

\newcommand{\YES}{\textsc{yes}}
\newcommand{\NO}{\textsc{no}}
\newcommand{\UNK}{\textsc{unknown}}
\newcommand{\Hb}{H_{\mathrm{b}}}

\title{Myopic Information-Gain Selection over a Theory Ensemble}
\author{}
\date{}

\begin{document}
\maketitle

\section*{Setup}

At a selection point we generate an ensemble of $M$ theories $T_1,\dots,T_M$,
ranked $1$ (most likely) to $M$. The candidate question set is the unanswered
bank $\mathcal{Q}$ (after it has been augmented with the crux questions).

\section{Prior over theories (rank $\to$ weight)}

\texttt{assign\_rank\_weights} puts a geometric prior on rank, with decay
$\gamma \in (0,1]$:
\[
\pi_k \;=\; P(T = T_k) \;=\; \frac{\gamma^{\,r_k - 1}}{\sum_{j=1}^{M}\gamma^{\,r_j - 1}},
\qquad r_k = \text{rank of } T_k .
\]
$\gamma < 1$ concentrates mass on top-ranked theories; $\gamma = 1$ gives a
uniform prior. (Default $\gamma = 0.6$.)

\section{Per-theory answer predictions}

For each candidate question $q_i$ and theory $T_k$, one
\texttt{\_predict\_answer\_vector} call yields a label in
$\{\YES, \NO, \UNK\}$, mapped to the theory's predicted probability that the
answer is \YES:
\[
a_{k,i} \;=\; P(A_i = \YES \mid T_k) \;=\;
\begin{cases}
1 & \YES \\[2pt]
0 & \NO \\[2pt]
\tfrac12 & \UNK \ (\text{agnostic})
\end{cases}
\]
So the per-theory binary answer model is $P(A_i{=}\YES \mid T_k) = a_{k,i}$,
$P(A_i{=}\NO \mid T_k) = 1 - a_{k,i}$. Cost: $M$ LLM calls total (independent
of $|\mathcal{Q}|$), versus B-diff's $\approx 4|\mathcal{Q}|$.

\section{Marginal answer and Bayesian posterior}

Marginal probability of a \YES{} under the prior:
\[
\rho_i \;=\; P(A_i = \YES) \;=\; \sum_{k=1}^{M}\pi_k\, a_{k,i}.
\]
Posterior over theories after observing each answer (Bayes):
\[
P(T_k \mid A_i{=}\YES) = \frac{\pi_k\, a_{k,i}}{\rho_i},
\qquad
P(T_k \mid A_i{=}\NO) = \frac{\pi_k\,(1 - a_{k,i})}{1 - \rho_i}.
\]

\section{Score = mutual information (expected info gain over theory identity)}

\[
\boxed{\;\text{score}(q_i) = I(A_i; T) = H(T) - \mathbb{E}_{A_i}\!\big[H(T \mid A_i)\big]\;}
\]
with everything in bits:
\[
H(T) = -\sum_k \pi_k \log_2 \pi_k,
\]
\[
H(T \mid A_i) = \rho_i\, H\!\big(T \mid A_i{=}\YES\big)
+ (1 - \rho_i)\, H\!\big(T \mid A_i{=}\NO\big),
\]
\[
H(T \mid A_i{=}a) = -\sum_k P(T_k \mid A_i{=}a)\,\log_2 P(T_k \mid A_i{=}a).
\]
This is exactly \texttt{mutual\_information(weights, p\_yes\_per\_theory)}: it
normalizes $\text{yes\_weights} = \pi_k a_{k,i}$ and
$\text{no\_weights} = \pi_k(1 - a_{k,i})$ to get the two posteriors, weights
their entropies by $\rho_i,\ 1-\rho_i$, and returns
$\max\big(0,\ H(T) - H(T \mid A_i)\big)$.

\paragraph{Equivalent dual form} (cheaper to state, useful as a sanity check ---
the code uses the form above):
\[
I(A_i; T) = H(A_i) - H(A_i \mid T)
= \Hb(\rho_i) \;-\; \sum_k \pi_k\, \Hb(a_{k,i}),
\]
where $\Hb(p) = -p\log_2 p - (1-p)\log_2(1-p)$ is the binary entropy. From this
it is immediate that:
\begin{itemize}[nosep]
  \item theories agree ($a_{k,i}$ all equal) $\Rightarrow H(A_i \mid T) = \Hb(\rho_i)
        \Rightarrow I = 0$ (\texttt{test\_mi\_zero\_when\_theories\_agree});
  \item two equal-weight theories that disagree ($a = \{1,0\}$,
        $\pi = \{\tfrac12,\tfrac12\}$) $\Rightarrow I = \Hb(\tfrac12) = 1$ bit
        (\texttt{test\_mi\_max\_for\_even\_split});
  \item an all-\UNK{} question scores $0$ (each $\Hb(\tfrac12)$ cancels).
\end{itemize}

\section{Selection}

Probe-selection (the LLM) narrows the bank to a candidate set $S$; the pursued
target is
\[
q^\star = \arg\max_{q_i \in S} I(A_i; T),
\]
with an LLM tie-break if the top score is shared. The experiment plan is then
formulated from $q^\star$. The crux questions are, by construction, drafted to
make $a_{\cdot,i}$ disagree across theories --- i.e.\ to be high-$I(A_i; T)$
candidates injected directly into $\mathcal{Q}$.

\section*{What this is not (the B vs.\ A boundary)}

This is the myopic, one-step, stateless version. The prior $\pi$ comes only
from the LLM's rank ordering, and $T$ is regenerated fresh each selection point
--- there is no carried-forward posterior update from the actually observed
outcome:
\[
\pi_k \;\not\leftarrow\; \frac{\pi_k\, P(o_t \mid T_k)}{\sum_j \pi_j\, P(o_t \mid T_j)}.
\]
Adding that recursive Bayesian update (theories persist across steps, weights
updated by observed transitions, action chosen to maximize $I$ over a horizon)
is precisely the step up to Plan A.

\end{document}
